\documentclass[12pt,a4paper]{article}

\usepackage{cmap}
\usepackage[T2A]{fontenc}
\usepackage[utf8]{inputenc}
\usepackage[russian]{babel}
\usepackage{amsmath,amssymb}
\usepackage[a4paper,margin=2cm]{geometry}

\setlength{\parindent}{0pt}
\setlength{\parskip}{0.45em}
\sloppy

\begin{document}

\begin{center}
{\Large\bfseries Билет 1: вопросы для проверки знаний}
\end{center}

\noindent\fbox{%
  \begin{minipage}{\dimexpr\linewidth-2\fboxsep-2\fboxrule\relax}
  \normalfont
Предел последовательности точек в $\mathbb{R}^n$. Связь между сходимостью
последовательности точек и сходимостью последовательностей их координат.
Внутренние, предельные, изолированные точки множества. Открытые и замкнутые
множества, их свойства. Непрерывная кривая, область, компакт. Критерий
компактности в $\mathbb{R}^n$. Внутренность, замыкание и граница множества.
  \end{minipage}%
}

\section*{Вопросы на понимание}

\begin{enumerate}
\item Чем предельная точка множества отличается от точки прикосновения?
\item Может ли точка прикосновения не принадлежать множеству? Приведите пример.
\item Может ли предельная точка не принадлежать множеству? Приведите пример.
\item Может ли изолированная точка быть предельной точкой того же множества?
\item Почему граничная точка множества не обязана принадлежать самому множеству?
\item Найдите внутренность, замыкание и границу интервала $(0,1)$ как подмножества $\mathbb{R}$.
\item Найдите внутренность, замыкание и границу множества $[0,1]$ как подмножества $\mathbb{R}$.
\item Найдите внутренность, замыкание и границу открытого шара в $\mathbb{R}^n$.
\item Является ли множество $(0,1)$ компактом в $\mathbb{R}$? Почему?
\item Является ли множество $[0,1]$ компактом в $\mathbb{R}$? Почему?
\item Является ли множество $\mathbb{R}^n$ замкнутым? открытым? компактным?
\item Почему в критерии компактности в $\mathbb{R}^n$ нужны оба условия: замкнутость и ограниченность?
\end{enumerate}

\section*{Источники для сверки}

\texttt{target/answer\_question\_01.tex};
\texttt{IvGE\_1.pdf}, стр. 144--145, 158--164, 173, 177, 335;
\texttt{IvGE\_2.pdf}, стр. 56--57.

\end{document}
